The \textbf{BLUE\_DB} dataset (\textit{Basic Local Unit Election Database}) provides
structured, unit-level electoral results for national legislative elections and
European Parliament elections in European countries, together with a harmonised
geographic reference frame based on EU GISCO identifiers.

\medskip
\noindent
The dataset covers \VAR{n_countries} countries from \VAR{year_min} to
\VAR{year_max}, comprising \VAR{n_national_elections} national elections and
\VAR{n_eu_elections} European Parliament elections.  Results are reported at
the level of municipalities (LAU~2) where available, or at the lowest
geographic level published by national statistical offices.  Special aggregate
units---such as postal votes, overseas voter counts, and electoral
constituencies---are included in a separate reference file and linked to the
main geographic hierarchy.

\medskip
\noindent
Geographic units are identified by \textit{GISCO IDs}, a harmonised coding
scheme derived from Eurostat's GISCO database.  A GISCO ID takes the form
\texttt{CC\_LLLLLL}, where \texttt{CC} is the ISO~3166-1 alpha-2 country code
and \texttt{LLLLLL} is a numeric or alphanumeric local unit code.  Special and
aggregate units use codes in the \texttt{CC\_CCZZXX} range following the
conventions described in Section~\ref{sec:special}.

\medskip
\noindent
Party results are identified by \texttt{party\_id}, a unique country-level
key that harmonises source-file aliases to one party entry in
\texttt{parties.csv}.  The parties file links each entry to Wikidata,
European party affiliation, and EP parliamentary-group metadata.

\medskip
\noindent
\textbf{Coverage note.}  Not all countries are covered for all years.
Country-specific notes in Chapter~\ref{ch:countries} document the elections
included, known data-quality issues, and deviations from the standard format.
